% Part: computability
% Chapter: computability-theory
% Section: total

\documentclass[../../../include/open-logic-section]{subfiles}

\begin{document}

\olfileid{cmp}{thy}{tot}
\olsection{عدم قابلية تقرير الكلية}

ومن استعمالات مبرهنة ${\OLId{latin-ordinary}{s}}$-${\OLId{latin-ordinary}{m}}$-${\OLId{latin-ordinary}{n}}$ في إثبات عدم قابلية
الحساب المثال الآتي. نجعل $\fn{Tot}$ مجموعة فهارس الدوال الكلية القابلة
للحساب، فيكون
\[
\fn{Tot} = \Setabs{{\OLId{latin-ordinary}{x}}}{\text{لكل ${\OLId{latin-ordinary}{y}}$، تكون $\cfind{{\OLId{latin-ordinary}{x}}}({\OLId{latin-ordinary}{y}})\fdefined$}}.
\]

\begin{prop}
\ollabel{prop:total}
المجموعة $\fn{Tot}$ غير قابلة للحساب.
\end{prop}

\begin{proof}
نثبت انتفاء قابلية $\fn{Tot}$ للحساب باختزال ${\OLId{latin-looped}{K}}$
إليها. فنعرّف ${\OLId{latin-ordinary}{h}}({\OLId{latin-ordinary}{x}},{\OLId{latin-ordinary}{y}})$ على الوجه الآتي:
\[
{\OLId{latin-ordinary}{h}}({\OLId{latin-ordinary}{x}},{\OLId{latin-ordinary}{y}}) \simeq
\begin{cases}
{\OLMathNumeral{0}} & \text{إذا كان ${\OLId{latin-ordinary}{x}} \in {\OLId{latin-looped}{K}}$} \\
\fundefined & \text{خلاف ذلك.}
\end{cases}
\]
أما ${\OLId{latin-ordinary}{h}}({\OLId{latin-ordinary}{x}},{\OLId{latin-ordinary}{y}})$ فلا اعتماد لها على ${\OLId{latin-ordinary}{y}}$ أصلًا.
ومن اليسير حساب ${\OLId{latin-ordinary}{h}}$ جزئيًا: نتلقى ${\OLId{latin-ordinary}{x}},{\OLId{latin-ordinary}{y}}$، ونحسب~${\OLId{latin-ordinary}{h}}$ بمحاكاة
الدالة~$\cfind{{\OLId{latin-ordinary}{x}}}$ عند الدخل~${\OLId{latin-ordinary}{x}}$؛ فإن توقف الحساب، أعطت
${\OLId{latin-ordinary}{h}}({\OLId{latin-ordinary}{x}},{\OLId{latin-ordinary}{y}})$ القيمة~${\OLMathNumeral{0}}$ ووقفت. فالدالة ${\OLId{latin-ordinary}{h}}({\OLId{latin-ordinary}{x}},{\OLId{latin-ordinary}{y}})$ هي
$\Zero(\umin{{\OLId{latin-ordinary}{s}}}{{\OLId{latin-looped}{T}}({\OLId{latin-ordinary}{x}},{\OLId{latin-ordinary}{x}},{\OLId{latin-ordinary}{s}})})$، و$\Zero$ دالة الصفر الثابتة.

ومن مبرهنة ${\OLId{latin-ordinary}{s}}$-${\OLId{latin-ordinary}{m}}$-${\OLId{latin-ordinary}{n}}$ نحصل على دالة عودية بدائية~${\OLId{latin-ordinary}{k}}({\OLId{latin-ordinary}{x}})$، بحيث لكل
${\OLId{latin-ordinary}{x}}$ و~${\OLId{latin-ordinary}{y}}$،
\[
\cfind{{\OLId{latin-ordinary}{k}}({\OLId{latin-ordinary}{x}})}({\OLId{latin-ordinary}{y}}) \simeq
\begin{cases}
{\OLMathNumeral{0}} & \text{إذا كان ${\OLId{latin-ordinary}{x}} \in {\OLId{latin-looped}{K}}$} \\
\fundefined & \text{خلاف ذلك.}
\end{cases}
\]
فتكون $\cfind{{\OLId{latin-ordinary}{k}}({\OLId{latin-ordinary}{x}})}$ كلية حين ${\OLId{latin-ordinary}{x}} \in {\OLId{latin-looped}{K}}$، وغير معرّفة في أي موضع
إن لم يكن. فهذه ${\OLId{latin-ordinary}{k}}$ اختزال من ${\OLId{latin-looped}{K}}$ إلى~$\fn{Tot}$.
\end{proof}

\begin{digress}
بل إن $\fn{Tot}$ لا تقبل التعداد الحسابي أصلًا؛ فتعقيدها في
منزلة أعلى من «الهرم الحسابي». ولا نتناول هنا إثبات هذه الزيادة.
\end{digress}

\end{document}
